\thispagestyle{plain}

\begin{center}
    \large
    \textbf{Feladatkiírás}
    \phantomsection{}
    \addcontentsline{toc}{chapter}{Feladatkiírás}
\end{center}

A feladat egy elosztott teszt rendszer fejlesztése, ami képes fordítóprogramok tesztelésére. A tesztrendszer bemenete a
teszt forrásfájlok, egy \textquote{\textit{SET}} fájl, ami meghatározza az egy teszt csomagba tartozó forrásfájlokat,
valamint egy konfigurációs fájl, ami megszorításokat, kombinációkat és kapcsolókat ad a teszt csomag futtatására. A
rendszer képes a különböző forrás fájlokat és kapcsolókat átadni a tesztelendő fordítóprogram számára, majd ellenőrizni
a fordítóprogram kimenetét és kategorizálni azt. A rendszer három különböző teszt típust képes kezelni:

\begin{itemize}
    \item Fordítási teszt, elvárt viselkedés, hogy tárgykód jön létre.
    \item Fordítási és linkelési teszt, elvárt viselkedés, hogy futtatható bináris jön létre.
    \item Futtatható teszt, elvárt viselkedés, hogy a lefordított tesztet valamilyen futás idejű logika ellenőrzi.
\end{itemize}

A futás idejű tesztek egy vékony \textquote{\textit{runtime}} könyvtáron keresztül működnek, míg a fordítási és
linkelési idejű teszteknél vagy diagnosztikai üzenet, vagy hiba nélküli futás az elvárt viselkedés.

A harmadik pontnak megfelelően a rendszer képes egy megfelelő emulátor használatára, ha a futtatható állomány nem a
futtató platformra készült. Ezenfelül a rendszer képes véletlenszerűen generált tesztek generálására a
\textquote{\textit{Csmith}} programon keresztül. A rendszer egy egyszerű kliens-szerver protokollon keresztül képes
elosztottan teszteket futtatni.

\begin{center}
    \large
    \textbf{Tartalmi összefoglaló}
    \phantomsection{}
    \addcontentsline{toc}{chapter}{Tartalmi összefoglaló}
\end{center}

A szakdolgozat témája egy elosztott fordítóprogram teszt rendszer fejlesztésével foglalkozik. A feladat megismerése
után először az alap funkcionalitás, a tesztek futtatását fogom implementálni. Erre az alapra építhető a konfigurációs
fájl megadása, a tesztek kategorizálása, a véletlenszerűen generált tesztek előállítása, és az alkalmazás elosztott
kliens-szerver részének fejlesztése.

A fejlesztés Fedora 41 operációs rendszeren folyik, Python 3.12 nyelven, a nyelv viszonylag új statikus típusosságát
használva. A fejlesztéshez Visual Studio Code fejlesztő környezetet használok, python programozáshoz specifikus
bővítményekkel. Debuggoláshoz a debugpy eszközt használom.

\begin{itemize}
    \item \textit{\textbf{Kulcsszavak:} fordítóprogram, elosztott alkalmazás, teszt rendszer, statikus python}
\end{itemize}
